%%%%%%%%%%%%%%%%%%%%%%%%%%%%%%%%%%%%%%%%%%%%%%%%%%%%%%%%%%%%
%% CHAPTER 3 -- NEURAL NETWORKS
%% Working draft. Not yet included in main.tex; will be
%% integrated, with renumbering of existing CH3-CH12 to
%% CH4-CH13, in Task 12 of the Wave 11 task list.
%%
%% Currently contains: Setting the Scene prologue (Task 4).
%% Body sections, scaffolding, worked examples, and end matter
%% will be appended in Tasks 5-11.
%%%%%%%%%%%%%%%%%%%%%%%%%%%%%%%%%%%%%%%%%%%%%%%%%%%%%%%%%%%%

\section{Neural Networks}
\label{sec:neural-networks}

\subsection*{Setting the Scene}
\addcontentsline{toc}{subsection}{Setting the Scene}

\paragraph{The problem.}
You have, by this point, two of the three things you need to train
a Large Language Model. From Chapter~1 you have the loss --- the
cross-entropy that asks, of any probability distribution over next
tokens, how surprised it was by the token that actually came next.
From Chapter~2 you have the optimiser --- stochastic gradient
descent, with its learning-rate schedules and its bias-corrected
moments and its Chinchilla-shaped compute budgets. What you do not
yet have is the \emph{function class} that the optimiser is
updating. That is the subject of this chapter. Before we can
compose the Transformer in Chapter~4, we need to introduce, one at
a time, the pieces it composes: a layer, a stack of layers, a way
to compute gradients through that stack, a choice of activation, a
method for initialising the weights, a normalisation, and a trick
for keeping deep stacks trainable. This chapter is the place where
each of those pieces enters the book.

\paragraph{How we got here.}
The idea that a network of simple units could learn a function from
examples is older than the computer that could run it. Frank
Rosenblatt, at Cornell in 1958, built the perceptron --- a single
linear threshold unit that could be trained to classify inputs by
adjusting its weights in proportion to its errors. The perceptron
was, briefly, a sensation. It was also limited: in 1969 Marvin
Minsky and Seymour Papert proved that no single linear threshold
unit could represent the exclusive-or of two inputs, and in the
backlash that followed the field receded for a decade.

The decisive response was depth. Paul Werbos, in his 1974 Harvard
PhD thesis, wrote down a recognisable derivation of what we now
call backpropagation --- the algorithm by which the chain rule of
calculus computes gradients through a multi-layer composition.
Backpropagation sat largely unread until 1986, when David
Rumelhart, Geoffrey Hinton, and Ronald Williams published a paper
in \emph{Nature} that brought it to the attention of an audience
ready to use it. Yann LeCun, at AT\&T, showed in 1989 that the
same algorithm could train a network to read handwritten zip
codes; the field had its first large-scale demonstration that
depth on real problems was tractable. Sepp Hochreiter and J\"urgen
Schmidhuber, in 1997, diagnosed the vanishing-gradient problem in
deep recurrent networks and proposed the gated cell that solved
it.

The next move was to make plain feed-forward depth work without
elaborate pre-training. Geoffrey Hinton, Simon Osindero, and
Yee-Whye Teh proposed a layer-wise pre-training recipe in 2006
that briefly closed the gap. Xavier Glorot and Yoshua Bengio in
2010 introduced the variance-preserving initialisation that made
the recipe unnecessary; the same year, Vinod Nair and Geoffrey
Hinton showed that the rectified linear unit --- a function so
simple it had been overlooked --- removed one of the main
bottlenecks in the gradient signal. The 2012 AlexNet paper, by
Alex Krizhevsky, Ilya Sutskever, and Geoffrey Hinton, took these
pieces to ImageNet at scale and won by a margin that ended the
debate about whether deep networks worked.

The remaining pieces converged in the years that followed. Sergey
Ioffe and Christian Szegedy added batch normalisation in 2015.
Kaiming He and his coauthors gave us, in two papers that same year,
both the ReLU-tuned variant of the Glorot init and the residual
connection that made networks of more than a hundred layers
trainable. Jimmy Ba, Jamie Ryan Kiros, and Geoffrey Hinton in 2016
introduced layer normalisation, the variant that survived into
the Transformer. Dan Hendrycks and Kevin Gimpel, also in 2016,
replaced the rectifier with a smoothed version called GeLU. Biao
Zhang and Rico Sennrich, in 2019, simplified layer normalisation
into RMSNorm. And Noam Shazeer, in 2020, identified the gated
activation variants (SwiGLU among them) that now sit inside almost
every Transformer's feed-forward block. Six decades of named
contributions, and the modern neural network is what you get when
you compose them.

\paragraph{Where we stand.}
A modern neural network used inside a Large Language Model is a
stack of differentiable transformations. Each transformation is a
linear layer (multiply by a weight matrix, add a bias vector)
followed by an elementwise nonlinearity. The whole network is
trained by gradient descent, with the gradients computed by
backpropagation. The depths that turned out to be useful --- tens
to hundreds of layers --- are made trainable by three engineering
moves: weight initialisation that preserves the variance of
activations across layers, a normalisation applied at each layer
to re-establish that variance after training has perturbed it, and
a residual connection that lets the optimiser learn small
corrections to an identity mapping rather than the full
transformation outright. Each of these moves was a response to a
specific failure of the design that came before; none was
inevitable. The Transformer of Chapter~4 is what you get when you
take this neural-network template and add the attention mechanism
on top.

\paragraph{What goes wrong.}
The failure modes of neural networks are the engineering
vocabulary of the field, and they recur often enough that you
should learn to recognise them. The first is \emph{vanishing
gradients}: in a deep network with the wrong initialisation or
the wrong activation, the gradient signal shrinks exponentially
with depth, and the early layers cease to learn. The second is
\emph{exploding gradients}: the same mechanism in reverse, with
activations and gradients growing past the range of floating-point
arithmetic until the network produces NaNs and the run dies. The
third is the \emph{dead-neuron problem}: a unit whose
pre-activation is permanently negative under a ReLU receives zero
gradient forever, becoming an ineffective contributor. The fourth
is the \emph{expressive-versus-trainable} gap: a network
architecture that \emph{can} represent the function we want is
not necessarily a network that the optimiser \emph{will find}
such a representation in. Each of these failures has a concrete
remedy --- careful initialisation, the right activation,
normalisation, residual paths --- and each remedy is one of the
moves this chapter unpacks.

Before that unpacking begins, a note on prerequisites. If you
have taken a course in deep learning and can write down
backpropagation on a whiteboard, you can skim or skip this
chapter; the material that follows is the foundational floor the
rest of the book stands on, and it will not surprise you. If you
have not, this is the chapter where the rest of the book becomes
legible. Either way, the goal is the same: to arrive at Chapter~4
with the function class in hand, ready to meet the architecture
that has, for now, eaten the field.

% --- END Setting the Scene ---
% Subsequent sections (Learning Objectives, Components Covered,
% Notation, §1-§7, Worked Examples, end matter) follow in
% Tasks 5-11.
